% Part: first-order-logic
% Chapter: proof-systems
% Section: tableaux

\documentclass[../../../include/open-logic-section]{subfiles}

\begin{document}

\iftag{FOL}
      {\olfileid{fol}{prf}{tab}}
      {\olfileid{pl}{prf}{tab}}

\olsection{\usetoken{P}{tableau}}

ઘણાં !!{derivation} તંત્રો !!{sentence}sની ગોઠવણી પર કાર્ય કરે છે,
જ્યારે !!{tableau}s !!{signed formula}s પર કાર્ય કરે છે.
!!^a{signed formula} એ સત્યમૂલ્યના ચિહ્ન ($\True$ અથવા $\False$) અને
!!a{sentence}ની બનેલી જોડી છે:
\[
\sFmla{\True}{!A} \text{ અથવા } \sFmla{\False}{!A}.
\]
!!^a{tableau} નીચે તરફ શાખાઓ પાડતા વૃક્ષમાં ગોઠવેલાં
!!{signed formula}sનું બનેલું હોય છે. તે કેટલીક \emph{ધારણાઓ}થી શરૂ
થાય છે અને એવાં !!{signed formula}s સાથે આગળ વધે છે જે ઉપરનાં
!!{signed formula}sમાંથી કોઈ એક પર અનુમાનનો કોઈ નિયમ લગાડવાથી મળે છે.
દરેક નિયમથી શાખાના છેડે એક કે વધુ
!!{signed formula}s, અથવા બાજુબાજુ બે !!{signed formula}s ઉમેરી શકાય
છે. પછીના કિસ્સામાં શાખા બે ભાગમાં વહેંચાય છે અને ઉમેરેલાં બે !!{signed
  formula}s બંને શાખાના છેડા બને છે.

સંકુલ !!{signed formula} પર લગાડેલો નિયમ એવાં !!{signed formula}s
ઉમેરે છે જે તેનાં સીધાં ઉપ-!!{formula}s હોય. નિયમો જોડમાં આવે છે—બે
ચિહ્નોમાંથી દરેક માટે એક નિયમ. દાખલા તરીકે, નિયમ
$\TRule{\True}{\land}$ એ $\sFmla{\True}{!A \land !B}$ને લાગુ પડે છે,
અને $\sFmla{\True}{!A \land !B}$ ધરાવતી કોઈપણ શાખાના છેડે બંને
!!{signed formula}s $\sFmla{\True}{!A}$ અને~$\sFmla{\True}{!B}$
ઉમેરવાની છૂટ આપે છે. નિયમ $\TRule{\False}{\land}$ શાખાને બે
ભાગમાં વહેંચીને બાજુબાજુ $\sFmla{\False}{!A}$ અને
$\sFmla{\False}{!B}$ ઉમેરવાની છૂટ આપે છે. જો !!^a{tableau}ની દરેક
શાખામાં !!{signed formula}sની મેળ ખાતી જોડી
$\sFmla{\True}{!A}$ અને $\sFmla{\False}{!A}$ હોય, તો તે ટેબ્લો બંધ છે.
\sourcecorrection{OLPRF-002}{મૂળની \texttt{tableaux.tex}, પંક્તિ~38માં
અસત્ય ચિહ્નવાળા સંયોજનના નિયમનું નામ $\TRule{\False}{!A\land !B}$
લખાયું છે. નિયમોની નિયામક યાદી અને આ જ ગ્રંથના બધા ઉપયોગોમાં નિયમનું
નામ $\TRule{\False}{\land}$ છે; અહીં તે જ નામ પુનઃસ્થાપિત કર્યું છે.}

!!{tableau}s પર આધારિત $\Proves$ સંબંધ આ પ્રમાણે વ્યાખ્યાયિત થાય છે:
$\Gamma \Proves !A$ જો અને માત્ર જો એવો કોઈ સાન્ત ગણ
$\Gamma_0 = \{!B_1, \dots, !B_n\} \subseteq \Gamma$ હોય કે ધારણાઓ
\[
\{\sFmla{\False}{!A}, \sFmla{\True}{!B_1}, \dots, \sFmla{\True}{!B_n}\}
\]
માટે કોઈ બંધ !!{tableau} હોય. દાખલા તરીકે, નીચેનો બંધ !!{tableau}
બતાવે છે કે $\Proves (!A \land !B) \lif !A$:
\begin{oltableau}
  [\sFmla{\False}{(\formula{A} \land \formula{B}) \lif \formula{A}}, just = \TAss
    [\sFmla{\True}{\formula{A} \land \formula{B}}, just = {\TRule{\False}{\lif}[1]}
      [\sFmla{\False}{\formula{A}}, just = {\TRule{\False}{\lif}[1]}
        [\sFmla{\True}{\formula{A}}, just = {\TRule{\True}{\land}[2]}
          [\sFmla{\True}{\formula{B}}, just = {\TRule{\True}{\land}[2]}, close]
        ]
      ]
    ]
  ]
\end{oltableau}
\sourcecorrection{OLPRF-004}{મૂળની \texttt{tableaux.tex},
પંક્તિઓ~59--60માં બીજી પંક્તિના સત્ય-સંયોજનને વિઘટિત કરતી બંને
શાખાઓ પર $\TRule{\True}{\lif}[2]$ લખાયું છે. અહીં વપરાયેલું નિયમ
$\TRule{\True}{\land}[2]$ છે; બંને નિયમ-નામ તે મુજબ સુધાર્યાં છે.}

ધારણાઓ
\[
\{\sFmla{\True}{!B_1}, \dots, \sFmla{\True}{!B_n}\}
\]
માટે બંધ !!{tableau} હોય અને દરેક $i=1,\dots,n$ માટે
$!B_i\in\Gamma$ હોય, તો ગણ $\Gamma$ !!{tableau} કલનમાં અસુસંગત છે.
\sourcecorrection{OLPRF-003}{મૂળની \texttt{tableaux.tex}, પંક્તિ~70માં
ઉપર દર્શાવેલી બધી ધારણાઓમાંથી માત્ર ``કોઈ એક $!B_i\in\Gamma$'' હોવાની
શરત લખાઈ છે. આ જ ગ્રંથની નિયામક સુસંગતતા-વ્યાખ્યા મુજબ સાન્ત ગણ
$\{!B_1,\dots,!B_n\}$ આખો $\Gamma$નો ઉપગણ હોવો જોઈએ; તેથી અહીં
દરેક $i=1,\dots,n$ માટે સભ્યતાની શરત લખી છે.}

!!^{tableau}sની શોધ 1950ના દાયકામાં એવર્ટ બેથ અને યાક્કો હિન્ટિકાએ
સ્વતંત્ર રીતે કરી હતી; રેમન્ડ સ્મલ્યને તેમને સરળ બનાવી લોકપ્રિય કર્યા.
!!a{tableau} રચવાની પ્રક્રિયા બહુ પદ્ધતિસર હોવાથી તેનો ઉપયોગ ઘણો સહેલો
છે. આ પદ્ધતિસર સ્વરૂપને કારણે !!{tableau}sનો કમ્પ્યૂટર દ્વારા અમલ પણ
સહેલાઈથી થઈ શકે છે. પરંતુ !!a{tableau} વાંચવો ઘણી વાર અઘરો હોય છે અને
સાબિતીઓ સાથેનો તેનો સંબંધ ક્યારેક સરળતાથી દેખાતો નથી. આ અભિગમ ઘણો
વ્યાપક પણ છે, અને ઘણાં જુદાં તર્કશાસ્ત્રોને !!{tableau} તંત્રો છે.
!!^{tableau}s એવી !!{structure}s શોધવામાં પણ મદદ કરે છે જે આપેલાં
!!{sentence}sના (ગણોને) સંતોષે છે: ગણ સંતોષ્ય હોય તો તેને બંધ !!{tableau} નહિ
હોય, એટલે કે દરેક !!{tableau}માં કોઈ ખુલ્લી શાખા હશે. તે શાખા પર
લગાડી શકાય એવો દરેક નિયમ લગાડ્યો હોય, તો ખુલ્લી શાખામાંથી સંતોષતી
!!{structure} ``વાંચી કાઢી'' શકાય છે. સિક્વન્ટ કલન સાથે પણ બહુ નજીકનો
સંબંધ છે: મૂળભૂત રીતે, બંધ !!{tableau} એ ઊંધું લખેલું, સંક્ષિપ્ત કરેલું
સિક્વન્ટ-કલનનું !!{derivation} છે.

\end{document}
